\section{Methodik}
\label{sec:methodik}

Nach der gefunden Startlösung wird die Tabu-Search-Heuristik angewendet, um eine verbesserte Lösung zu finden. Der Algorithmus ist in Algorithmus \ref{alg:tabu_search} dargestellt. Die Notation der Tabu-Search-Heuristik ist in Tabelle \ref{tab:notation_tabu} zusammengefasst.

\begin{algorithm}[H]
\caption{Tabu Search für das Job-Shop-Scheduling-Problem mit sequenzabhängigen Rüstzeiten}
\label{alg:tabu_search}

\begin{algorithmic}[1]

\State $S \gets \hat{S}$
\State $S^\ast \gets S$
\State $TL \gets \emptyset$
\State $I_{\text{total}} \gets 0$
\State $I \gets 0$
\State $\tau \gets 0$

\While{$I_{\text{total}} < I_{\text{total}}^{\max}$ \textbf{and} $I < I^{\max}$ \textbf{and} $\tau < \tau^{\max}$}
    \State $I_{\text{total}} \gets I_{\text{total}} + 1$
    \State Erzeuge die Nachbarschaft $NB(S)$ der aktuellen Lösung

    \If{$I \geq I^{\text{div}}$}
        \State Ergänze zufällige nicht-kritische Bewegungen zu $NB(S)$
    \EndIf

    \For{$S' \in NB(S)$} (jede Nachbarlösung)
        \If{$S'$ ist unzulässig \textbf{ or } (tabuiert \textbf{ and } Aspiration nicht erfüllt)}
            \State Entferne $S'$ aus $NB(S)$
        \EndIf
    \EndFor

    \State Wähle die beste zulässige Nachbarlösung $S' \in NB(S)$

    \State $S \gets S'$
    \State Entferne abgelaufene Einträge aus $TL$
    \State $TL \gets TL \cup \{\text{inverse Bewegung von } S'\}$

    \If{$F(S) < F(S^\ast)$}
        \State $S^\ast \gets S$
        \State $I \gets 0$
    \Else
        \State $I \gets I + 1$
    \EndIf
\EndWhile

\State \Return $S^\ast$

\end{algorithmic}
\end{algorithm}

Die Suche beginnt mit der Erzeugung einer zulässigen Startlösung . Anschließend wird die beste gefundene Lösung mit  initialisiert. Danach werden die Tabu-Liste (TL), der Gesamtiterationszähler und der Zähler für Iterationen ohne Verbesserung (I) initialisiert. Die Suche wird fortgeführt, bis ein Abbruchkriterium erfüllt ist. Der Gesamtiterationszähler wird in jedem Schleifendurchlauf um eins erhöht.
 
 Der Algorithmus verwendet drei Abbruchkriterien. Erstens endet die Suche nach Erreichen der maximalen Gesamtzahl an Iterationen. Zweitens endet sie nach einer festgelegten Anzahl an Iterationen ohne Verbesserung. Drittens kann ein Zeitlimit gesetzt werden. Der Timer wird nur aktiviert, wenn beim Aufruf des Algorithmus ein Zeitlimit übergeben wird.
 
 In jeder Iteration wird zunächst die Nachbarschaft (NB(S)) der aktuellen Lösung erzeugt. Ergänzend werden potenzielle Kandidatenbewegungen gebildet. Wird über mehrere Iterationen keine Verbesserung erzielt, werden zusätzliche zufällige nicht-kritische Bewegungen ergänzt. Diese dienen der Diversifikation der Suche. Liegen mehrere zulässige Kandidaten mit identischem Makespan vor, werden diese zufällig durchmischt. Dadurch soll ein Festlaufen in lokalen Optima reduziert werden.
 
 Tabuierte Bewegungen werden verworfen, sofern die Aspirationsbedingung nicht erfüllt ist. Diese ist erfüllt, wenn eine tabuierte Bewegung zu einer Verbesserung der besten bisher gefundenen Lösung führt. Aus den verbleibenden Kandidaten wird die beste zulässige Nachbarlösung (S') ausgewählt. Existiert keine zulässige Nachbarlösung, wird die Suche beendet.
 
 Wird eine Nachbarlösung akzeptiert, wird die aktuelle Lösung durch (S') ersetzt. Anschließend wird die Tabu-Liste mit der inversen Bewegung aktualisiert. Dadurch wird ein direktes Zurückkehren in den vorherigen Zustand verhindert. Ist die neue Lösung besser als die bisher beste Lösung, gilt also , wird  durch (S) ersetzt und der Verbesserungszähler zurückgesetzt. Andernfalls wird der Verbesserungszähler erhöht. Nach Abschluss der Suche wird die beste gefundene Lösung (S) zurückgegeben.
